\section{Билет 32. Эффект Доплера. Формула частоты, когда приемник двигается вдоль одной прямой, а источник покоится (с выводом). Обобщенная формула, когда источник и приемник двигаются вдоль одной прямой.}

\begin{definition}
\textbf{Эффектом Доплера} называется изменение частоты волн, регистрируемых приёмником, вследствие относительного движения источника волн и приёмника. Явление открыто австрийским физиком К. Доплером в 1842 г.
\par\noindent\textbf{Пример:} При приближении поезда к наблюдателю звук его гудка кажется более высоким, а при удалении --- более низким, чем у неподвижного источника.
\end{definition}

\begin{theorem}[Движущийся приёмник, неподвижный источник]
Пусть источник звука покоится, а приёмник движется вдоль прямой, соединяющей их, со скоростью $v_{\text{пр}}$. Скорость звука в среде --- $v$, частота источника --- $\nu_0$, длина волны $\lambda_0 = v/\nu_0$.

\textbf{Случай 1: приёмник приближается к источнику.}
За время $t$ приёмник проходит расстояние $v_{\text{пр}} t$ навстречу волнам. Число волн, пересекающих приёмник:
\begin{equation}
    N = \frac{v t + v_{\text{пр}} t}{\lambda_0} = \frac{(v + v_{\text{пр}}) t}{\lambda_0}.
\end{equation}
Частота, регистрируемая приёмником:
\begin{equation}
    \nu = \frac{N}{t} = \frac{v + v_{\text{пр}}}{\lambda_0} = \nu_0 \frac{v + v_{\text{пр}}}{v}. \label{eq:doppler_approach}
\end{equation}

\textbf{Случай 2: приёмник удаляется от источника.}
Аналогично, но приёмник «догоняет» волны:
\begin{equation}
    \nu = \nu_0 \frac{v - v_{\text{пр}}}{v}. \label{eq:doppler_recede}
\end{equation}
\end{theorem}

\begin{theorem}[Обобщённая формула эффекта Доплера]
Пусть и источник, и приёмник движутся вдоль одной прямой со скоростями $v_{\text{ист}}$ и $v_{\text{пр}}$ соответственно. Положительное направление оси $x$ --- от источника к приёмнику.

\textbf{Вывод:}
\begin{enumerate}
    \item Длина волны, излучаемая движущимся источником:
    \begin{equation}
        \lambda = \frac{v \mp v_{\text{ист}}}{\nu_0},
    \end{equation}
    где знак «минус» --- при сближении (источник «догоняет» свои волны), «плюс» --- при удалении.
    \item Частота, регистрируемая движущимся приёмником:
    \begin{equation}
        \nu = \frac{v \pm v_{\text{пр}}}{\lambda}.
    \end{equation}
    Знак «плюс» --- при сближении, «минус» --- при удалении.
\end{enumerate}
Подставляя $\lambda$, получаем \textbf{общую формулу эффекта Доплера}:
\begin{equation}
    \nu = \nu_0 \frac{v \pm v_{\text{пр}}}{v \mp v_{\text{ист}}}. \label{eq:doppler_general}
\end{equation}
\end{theorem}

\begin{property}[Правило знаков]
В формуле \eqref{eq:doppler_general}:
\begin{itemize}
    \item В числителе: знак «$+$» --- если приёмник \textbf{движется к источнику}, «$-$» --- если \textbf{удаляется}.
    \item В знаменателе: знак «$-$» --- если источник \textbf{движется к приёмнику}, «$+$» --- если \textbf{удаляется}.
\end{itemize}
Если скорости направлены не вдоль соединяющей прямой, в формулу подставляются проекции скоростей на эту прямую: $v_{\text{пр},x}$, $v_{\text{ист},x}$.
\end{property}

\begin{remark}
\textbf{Применение эффекта Доплера:}
\begin{enumerate}
    \item \textbf{Радары и лидары} --- измерение скорости автомобилей, самолётов, атмосферных потоков.
    \item \textbf{Медицинская диагностика} --- ультразвуковая допплерография кровотока.
    \item \textbf{Астрофизика} --- определение скорости звёзд и галактик по смещению спектральных линий (красное/синее смещение).
    \item \textbf{Навигация} --- спутниковые системы (GPS/ГЛОНАСС) учитывают доплеровский сдвиг для точного позиционирования.
\end{enumerate}
\end{remark}

\begin{figure}[htbp]
    \centering
    \begin{tikzpicture}[
        scale=1.0,
        >=Stealth,
        font=\small,
        wave/.style={blue, thick},
        source/.style={red, fill=red, circle, inner sep=3pt},
        receiver/.style={black, fill=black, circle, inner sep=3pt},
        box/.style={draw, fill=yellow!10, rounded corners, inner sep=4pt, align=center}
    ]

    % ==========================================
    % ЧАСТЬ 1: ДВИЖУЩИЙСЯ ИСТОЧНИК
    % ==========================================
    \begin{scope}[shift={(0, 3.5)}]
        \node[font=\bfseries] at (0, 2.6) {1. Движущийся источник ($v_{\text{ист}}$)};

        % Волновые фронты (сжатие/растяжение)
        \draw[wave] (-0.6, 0) circle (0.6);
        \draw[wave] (-1.2, 0) circle (1.2);
        \draw[wave] (-1.8, 0) circle (1.8);
        \draw[wave] (-2.4, 0) circle (2.4);

        % Источник
        \fill[red] (0,0) circle (3pt);
        \node[above right=2pt] at (0,0) {$S$};

        % Скорость источника
        \draw[->, red, thick] (0.5, 0.3) -- (2.5, 0.3) node[midway, above] {$v_{\text{ист}}$};

        % Наблюдатели
        \fill[black] (2.8, 0) circle (3pt);
        \node[below] at (2.8, 0) {$P_1$};
        \node[red, font=\footnotesize] at (2.8, 0.8) {$\nu > \nu_0$};

        \fill[black] (-3.2, 0) circle (3pt);
        \node[below] at (-3.2, 0) {$P_2$};
        \node[blue, font=\footnotesize] at (-3.2, 0.8) {$\nu < \nu_0$};

        % Длины волн
        \draw[<->, green!60!black, thick] (0.6, 0.5) -- (1.2, 0.5) node[midway, above] {$\lambda'$};
        \draw[<->, green!60!black, thick] (-3.0, 0.5) -- (-2.4, 0.5) node[midway, above] {$\lambda''$};

        % Пояснение
        \node[align=center] at (0, -2.5) {Спереди: $\lambda$ сжата\\Сзади: $\lambda$ растянута};
    \end{scope}

    % ==========================================
    % ЧАСТЬ 2: ДВИЖУЩИЙСЯ ПРИЁМНИК
    % ==========================================
    \begin{scope}[shift={(0, -3.5)}]
        \node[font=\bfseries] at (0, 2.8) {2. Движущийся приёмник ($v_{\text{пр}}$)};

        % --- Случай А: Приближение (Левая сторона) ---
        \begin{scope}[shift={(-3.5, 0)}]
            \node[font=\bfseries] at (0, 2.2) {Приближение};

            % Источник (Неподвижен)
            \fill[red] (0,0) circle (3pt);
            \node[above] at (0,0) {$S$};

            % Волны
            \draw[wave] (0,0) circle (0.8);
            \draw[wave] (0,0) circle (1.4);
            \draw[wave] (0,0) circle (2.0);

            % Приёмник
            \fill[black] (3.5, 0) circle (3pt);
            \node[below] at (3.5, 0) {$P$};

            % Скорость волны (Синяя, вправо)
            \draw[blue, ->, thick] (1.5, 0.5) -- (2.5, 0.5) node[midway, above] {$v$};

            % Скорость приёмника (Красная, ВЛЕВО - навстречу)
            \draw[red, ->, thick] (3.5, -0.8) -- (2.5, -0.8)
                node[midway, below, font=\bfseries] {$v_{\text{пр}}$};

            % Формула
            \node[box] at (0, -1.8) {$\nu = \nu_0 \dfrac{v + v_{\text{пр}}}{v}$};
        \end{scope}

        % --- Случай Б: Удаление (Правая сторона) ---
        \begin{scope}[shift={(4.0, 0)}]
            \node[font=\bfseries] at (0, 2.2) {Удаление};

            % Источник
            \fill[red] (0,0) circle (3pt);
            \node[above] at (0,0) {$S$};

            % Волны
            \draw[wave] (0,0) circle (0.8);
            \draw[wave] (0,0) circle (1.4);
            \draw[wave] (0,0) circle (2.0);

            % Приёмник
            \fill[black] (3.5, 0) circle (3pt);
            \node[below] at (3.5, 0) {$P$};

            % Скорость волны (Синяя, вправо)
            \draw[blue, ->, thick] (1.5, 0.5) -- (2.5, 0.5) node[midway, above] {$v$};

            % Скорость приёмника (Красная, ВПРАВО - в ту же сторону)
            \draw[red, ->, thick] (3.5, -0.8) -- (4.5, -0.8)
                node[midway, below, font=\bfseries] {$v_{\text{пр}}$};

            % Формула
            \node[box] at (0, -1.8) {$\nu = \nu_0 \dfrac{v - v_{\text{пр}}}{v}$};
        \end{scope}

    \end{scope}
    \end{tikzpicture}
    \caption{Эффект Доплера. Верх: движение источника меняет длину волны. Низ: движение приёмника меняет относительную скорость встречи с фронтами.}
    \label{fig:doppler_effect}
\end{figure}

